\Annex{informative}{Implementation Documentation}{AnnexB}

\Sec{Implementation Limits}{AnnexB.Limits}

\p This section describes reasonable implementation limits that may be assumed
by a conforming implementation. These limits specify minimum required or maximum
allowed values for various implementation-defined parameters.

\p \ref{Basic.Types} declares a vector type as having required supported sizes
of 1 to 4 components. Implementations may support larger vector sizes but are
not required to do so. An implementation that supports larger vector sizes must
define its vector size limit.

\Sec{Observed Deviations}{AnnexB.Documentation}

\p This section provides documentation around implementation behaviors that are
known to differ between specific implementations. This includes breaking
changes between historical HLSL and standard HLSL where some compilers or
versions of compilers may not conform to standard HLSL, as well as observed
non-conformant behaviors in such compilers. This is not meant to be a
comprehensive list of bugs.

\Sub{Resource Objects}{AnnexB.Resources}

\p In violation of \ref{Basic.Types.Intangible}, DXC supports storing objects
that contain resources in \texttt{cbuffer} declarations. In this use case the
resource objects are not actually stored in the \texttt{cbuffer}, but instead
the resource is treated as a global resource declaration and name lookup
resolves the value inside the structure declaration to the implied global
resource declaration.
